Pelvic mechanical response reflects anatomical variation expressed through a specified load path. We examined sacroiliac joint responses in 278 imaging-derived pelves under five standardized loading configurations and in matched geometry--bone-density variants. Redistributing the same 800~N total acetabular force from bilateral to unilateral support amplified motion and asymmetry; transverse force pairs applied at different sites produced distinct component patterns. Explicit pelvic dimensions substantially attenuated age-adjusted sex-associated translation differences under localized loading, with all fully adjusted confidence intervals spanning zero. Preserving individual geometry while standardizing the bone-density field retained subject-specific translation closely (RMSE $\leq0.028$~mm; minimum identity-line $R^2=0.946$). In retrospective adjusted log--log analyses, all 20 sex-specific volume--motion slopes were negative after multiplicity correction. Bilateral standing-like slopes were numerically near simple fixed-force dimensional reference values, whereas several localized-load estimates were numerically steeper. These conditional relations describe size-dependent response under identical prescribed forces across subjects within each load case. Together, the comparisons support a description of pelvic mechanics organized by anatomical scale and individual architecture, with load path determining their mechanical expression. Sex-associated contrasts varied with anatomical adjustment and loading configuration. Interpretation remains conditional on passive static mechanics, the implemented density--modulus law and fixed cartilage, ligament and pretension parameters.
